\section{Computational Toolchain}

\subsection{Geometry and Meshing}

The geometry layer defines the physical domain of a computational problem and is responsible for producing the finite-element mesh used by the mechanics formulation. Geometry declarations are kept distinct from their numerical realization: the problem definition specifies the physical objects that constitute the model, while the geometry package determines how those objects are constructed, meshed, and converted into the finite-element representation required by DOLFINx.

A computational problem may contain multiple geometry objects. Each object is identified by a semantic name and a geometry type, together with the parameters required to define its physical dimensions and placement. For example, a cylindrical specimen and two rigid platens can be declared independently:

\begin{verbatim}
geometry:
  - name: specimen
    type: cylinder
    x0: [0.0, 0.0, 0.0]
    x1: [0.0, 0.0, 0.010]
    radius: 0.005

  - name: bottom_platen
    type: cylinder
    x0: [0.0, 0.0, -0.001]
    x1: [0.0, 0.0, 0.0]
    radius: 0.010

  - name: top_platen
    type: cylinder
    x0: [0.0, 0.0, 0.010]
    x1: [0.0, 0.0, 0.011]
    radius: 0.010
\end{verbatim}

The geometry package separates the declaration of these physical objects from mesh generation. The initial implementation is organized around the following responsibilities:

\begin{description}
    \item[\texttt{definitions.py}] Defines typed representations of geometry declarations. This layer describes what physical geometry has been requested and does not depend on Gmsh or DOLFINx.

    \item[\texttt{registry.py}] Maps geometry type names, such as \texttt{cylinder}, \texttt{stl}, or \texttt{obj}, to their corresponding geometry implementations.

    \item[\texttt{builders.py}] Provides the common realization interface used to construct geometric models from the typed geometry definitions.

    \item[\texttt{cylinder.py}] Implements analytic cylindrical geometry. A cylinder is defined parametrically by its two end points and radius.

    \item[\texttt{sources.py}] Provides interfaces for geometry originating from external surface or mesh sources, including formats such as STL and OBJ.

    \item[\texttt{meshing.py}] Applies the mesh discretization policy, generates the finite-element mesh, assigns semantic physical groups, and converts the resulting model to the DOLFINx mesh representation.
\end{description}

The resulting package therefore follows the separation

\begin{center}
\textit{physical geometry declaration}
$\;\longrightarrow\;$
\textit{geometry realization}
$\;\longrightarrow\;$
\textit{mesh generation}
$\;\longrightarrow\;$
\textit{DOLFINx mesh}.
\end{center}

The initial discretization strategy uses conforming three-dimensional tetrahedral elements with first-order geometry. This provides a straightforward baseline for irregular three-dimensional grain geometries while leaving higher-order and curved discretizations available for later development. Mesh resolution is specified independently from the physical geometry through global mesh defaults, with optional local overrides for individual objects or features.

Semantic physical groups are retained during mesh generation so that downstream boundary-condition, loading, and contact definitions can refer to physical entities by name rather than by Gmsh-specific integer identifiers. This provides a stable interface between the problem-definition layer and the mechanics formulation.

The geometry package ultimately hands the generated mesh and associated entity tags to DOLFINx. Function-space construction, variational forms, assembly, and solution therefore remain outside the geometry layer. Similarly, adaptive mesh refinement is not treated as a geometry operation; the geometry layer provides the initial discretization, while any solution-driven refinement belongs to the numerical solution workflow.

